\chapter{Study Area and Data}
\label{ch:data}

This chapter describes the study area, the spatial unit of analysis, and the data
that feed the pipeline. A recurring theme is data integrity: several of the results in
this thesis depend on corrections to the base geography and to specific indicators that
were discovered during a full reproducibility rebuild, and these are documented here so
that the numbers throughout are traceable.

\section{The Guadalajara Metropolitan Area}
\label{sec:area-zmg}

The primary study area is the Zona Metropolitana de Guadalajara (\zmg), Mexico's second
largest metropolitan area. Following the official delimitation, it comprises ten
municipalities of the state of Jalisco, spanning a dense central core (Guadalajara and
Zapopan) and a rapidly-urbanizing periphery (\cref{tab:munis}). Its mobility is marked
by high and rising car dependence, a large informally-operated concessioned bus system,
and, more recently, a small premium tier of bus rapid transit (Mi Macro) and light rail
(Mi Tren). This premium tier follows the wider Latin-American pattern in which bus rapid
transit has expanded rapidly as a cost-effective complement to rail
\parencite{hidalgo2013brt, cervero2013brt}; grafting it onto a largely concessioned base
network is precisely the setting that motivates a demand-driven rather than
supply-following approach to network planning.

\begin{table}[htbp]
\centering
\caption{The ten \zmg municipalities and their \ageb counts.}
\label{tab:munis}
\begin{tabular}{lS[table-format=3.0]}
\toprule
Municipality & {\ageb{}s} \\
\midrule
Zapopan                       & 481 \\
Guadalajara                   & 404 \\
Tlajomulco de Zúñiga          & 341 \\
San Pedro Tlaquepaque         & 213 \\
Tonalá                        & 207 \\
El Salto                      & 115 \\
Zapotlanejo                   & 41  \\
Ixtlahuacán de los Membrillos & 34  \\
Acatlán de Juárez             & 29  \\
Juanacatlán                   & 16  \\
\midrule
\textbf{Total}                & \textbf{1881} \\
\bottomrule
\end{tabular}
\end{table}

\section{The \ageb as unit of analysis}
\label{sec:area-ageb}

All indicators are aggregated to the \ageb (\emph{área geoestadística básica}), the
census enumeration area defined by INEGI. The \ageb is fine enough to resolve
intra-urban variation in demand and access, yet coarse enough that census attributes
are reported at its level. The analysis universe is restricted to urban \ageb{}s of the
ten \zmg municipalities: cells whose \code{CVE\_AGEB} carries an alphabetic suffix
(rural \ageb{}s) are excluded, as are all municipalities outside the metropolitan
delimitation. After this filter the universe is \num{1881} \ageb{}s
(\cref{tab:munis})---not the \num{2068} used in earlier drafts of this work, a
discrepancy explained in \cref{sec:area-integrity}.

\section{Data sources}
\label{sec:area-sources}

The pipeline deliberately separates \emph{Tier-1} data---census, economic, and
network variables available for any Mexican city---from the transit-supply data used
only to measure the coverage gap and to audit existing routes. \Cref{tab:data-sources}
lists the principal sources.

\begin{table}[htbp]
\centering
\caption{Principal data sources and their role in the pipeline.}
\label{tab:data-sources}
\begin{tabularx}{\linewidth}{llX}
\toprule
Source & Level & Role \\
\midrule
INEGI CPV2020    & \ageb / block & Population, youth share, vehicle ownership \\
DENUE            & establishment & Employment and POI attraction proxies \\
OpenStreetMap    & network       & Intersections, street density, drive graph \\
GTFS             & stop / route  & Existing accessibility and route audit \\
EOD 2022         & survey zone   & Gravity-model calibration \\
CONAPO / CONEVAL & \ageb         & Marginalization and social-lag (equity) \\
\bottomrule
\end{tabularx}
\end{table}

\begin{description}[leftmargin=1.4em, style=nextline]
  \item[Census and socio-demographics (INEGI CPV2020).] The 2020 population and
    housing census \parencite{inegi_cpv2020} supplies population, the youth share used to
    amplify trip generation, and household vehicle ownership (\code{VPH\_AUTOM} over
    \code{VIVPAR\_HAB}) used to derive transit propensity.
  \item[Economic activity (DENUE).] The national economic-unit directory
    \parencite{inegi_denue} provides establishment locations, aggregated to \ageb
    employment and point-of-interest and retail densities that drive trip attraction and
    accessibility opportunities.
  \item[Street network (OpenStreetMap).] The drive graph and derived Node indicators
    (intersection and street density) come from \textsc{osm} \parencite{openstreetmap},
    extracted for the ten municipalities with \textsc{osmnx} \parencite{boeing2017osmnx}.
  \item[Transit supply (GTFS).] The SITEUR \textsc{gtfs} feed defines the existing
    network used for the accessibility surface (\cref{sec:w3}) and the route audit
    (\cref{sec:w7}).
  \item[Origin--Destination survey (EOD 2022).] The 2022 metropolitan travel survey
    \parencite{imeplan_eod2022} provides observed zonal flows for calibrating the
    gravity model's distance decay (\cref{sec:w2}).
  \item[Equity indices (CONAPO, CONEVAL).] The marginalization index
    \parencite{conapo_marginacion2020} and the social-lag index
    \parencite{coneval_rezago2020} form the equity term (\cref{eq:equity}).
\end{description}

\section{Spatial conventions and data integrity}
\label{sec:area-integrity}

All computation uses the canonical projected reference system EPSG:6372 (Mexico
conic-equidistant); geographic coordinates (EPSG:4326) are used only for ingestion.
All geometry columns are spatially indexed.

A full from-scratch reproducibility rebuild surfaced three integrity defects that had
been silently present in earlier drafts, all now corrected. First, a municipality-code
typo in the base-geography filter had dropped an entire municipality on every rebuild;
fixing it, together with correctly excluding alphabetic-suffix \ageb{}s, gives the
\num{1881}-\ageb universe used throughout (earlier numbers computed against a
never-reproducible \num{2068}-row table should be read as superseded). Second, the
digital-elevation raster had been loaded under the wrong spatial reference, producing
geometrically meaningless slope values; the affected legacy model was retired rather
than re-run. Third, the marginalization index entered the equity term with an inverted
sign and an out-of-range zero-fill for unmatched \ageb{}s, which together had let the
equity bonus reward the \emph{least}-marginalized areas; correcting the direction and
using median imputation restores a monotonic, correctly-signed equity term. These
corrections materially affect several headline numbers and are the reason the results
in this thesis supersede any earlier figures for the same quantities.
